% Einheit 4 von 5 – Grundwert berechnen (bank/prozentrechnung/e4.jsonl, zusammenbau v0.1)
%% TODO Zweigzeile Teil 1: was hier gelernt wird (Satz in Schülersprache aus dem Einheitstitel) – Einheit 4
\einheitenkopf[e4]{Einheit 4 von 5 · Grundwert berechnen}
\zweigzeile{neu in diesem Jahr · P10 oft}
\setcounter{aufgabe}{19}

%% TODO Titel: Ich-kann-Satz fehlt in der Bank (Platzhalter: Kettenname) – Nr. 20
%% TODO Anweisung: ein Satz über der ersten Teilaufgabe fehlt in der Bank – Nr. 20
\begin{aufgabe}{Was ist gegeben, was gesucht?}
\begin{teile}
\teil $30\,\%$ der Klasse sind $9$ Kinder. Wie viele Kinder hat die Klasse? – Was ist gesucht?\\ \kreuz{Prozentwert}\\ \kreuz{Prozentsatz}\\ \kreuz{Grundwert}
\end{teile}
\end{aufgabe}

%% TODO Titel: Ich-kann-Satz fehlt in der Bank (Platzhalter: Kettenname) – Nr. 21
%% TODO Anweisung: ein Satz über der ersten Teilaufgabe fehlt in der Bank – Nr. 21
\begin{aufgabe}{Grundwert}
\swa{Grau sind $30\,\%$, das sind $6$ Kästchen – wie viele Kästchen hat der ganze Streifen? \leerfeld[Kästchen]}{\streifen[0]{30}{$0\,\%$}{$100\,\%$}}
%% TODO Darstellung für schwach fehlt in der Bank (prozentrechnung-e4-k2-s1-v1; Abschnitt „Für schwache Schüler“)
\swz{$50\,\%$ sind $35$ kg – wie viel ist das Ganze?}{}
%% TODO Darstellung für schwach fehlt in der Bank (prozentrechnung-e4-k2-s1-v2; Abschnitt „Für schwache Schüler“)
\swz{$25\,\%$ sind $12\,€$ – wie viel ist das Ganze?}{}
%% TODO Darstellung für schwach fehlt in der Bank (prozentrechnung-e4-k2-s1-v3; Abschnitt „Für schwache Schüler“)
\swz{$20\,\%$ sind $9$ m – wie viel ist das Ganze?}{}
%% TODO Darstellung für schwach fehlt in der Bank (prozentrechnung-e4-k2-s1-v4; Abschnitt „Für schwache Schüler“)
\swz{$10\,\%$ sind $16$ l – wie viel ist das Ganze?}{}
%% TODO Darstellung für schwach fehlt in der Bank (prozentrechnung-e4-k2-s1-v5; Abschnitt „Für schwache Schüler“)
\swz{$25\,\%$ sind $30$ Kinder – wie viele sind es insgesamt?}{}
\swfrage{Was bleibt gleich, was ändert sich?}
%% TODO Darstellung für schwach fehlt in der Bank (prozentrechnung-e4-k2-s2-v1; Abschnitt „Für schwache Schüler“)
\swz{$3\,\%$ sind $12$ kg – wie viel ist das Ganze?}{}
%% TODO Darstellung für schwach fehlt in der Bank (prozentrechnung-e4-k2-s3-v1; Abschnitt „Für schwache Schüler“)
\swz{$36\,\%$ sind $81$ kg – wie viel ist das Ganze?}{}
%% TODO Darstellung für schwach fehlt in der Bank (prozentrechnung-e4-k2-s4-v1; Abschnitt „Für schwache Schüler“)
\swz{$21$ Kinder fahren mit dem Rad, das sind $75\,\%$ der Klasse – wie viele Kinder hat die Klasse?}{}
%% TODO Darstellung für schwach fehlt in der Bank (prozentrechnung-e4-k2-s5-v1; Abschnitt „Für schwache Schüler“)
\swz{$40\,\%$ von $65$ Kindern gehen in die AG – wie viele?}{}
\swz[3]{Beim Kauf einer Mütze spart Mira $4\,€$, das sind $25\,\%$ des alten Preises. Wie viel kostete die Mütze vorher? \hfill (P10 2025 OS)}{}
\end{aufgabe}

%% TODO Titel: Ich-kann-Satz fehlt in der Bank (Platzhalter: Kettenname) – Nr. 22
%% TODO Anweisung: ein Satz über der ersten Teilaufgabe fehlt in der Bank – Nr. 22
\begin{aufgabe}{Grundwert – Fehler finden, Begründen, Darstellungswechsel, Anwendung}
\swz[3]{„$30\,\%$ sind $12$ kg – wie viel ist das Ganze?“ Sven: \rechnung{0{,}3 \cdot 12 &= 3{,}6\text{ kg}} Fehler? Wie heißt es richtig?}{}
\swfrage{Warum rechnet man beim Grundwert am Ende mal $100$?}
\swa{Grau sind $40\,\%$, das sind $18$ kg. Trage die Werte an den Streifen, lies das Ganze ab. \leerfeld[kg]}{\streifen{40}{$0$ kg}{?}}
\swz[3]{Tims Handy zeigt $35\,\%$ Akku; laut Anzeige reicht das noch für $7$ Stunden Musik. Wie lange hält ein voller Akku? Reicht er für eine Busfahrt von $18$ Stunden?}{}
\end{aufgabe}

\uebersichtskasten{Grundwert: erst 1 \% (Prozentwert : p), dann mal 100. \\ 36 € sind 45 \%: 1 \% = 0,80 € → 100 \% = 80 € \quad Streifen: 45 \% ≙ 36 €, 100 \% ≙ ? \\ Erst zuordnen: Was ist gegeben (W, p, G), was gesucht? Dann rechnen.}
